\section{Limits and outlook}\label{sec: limits}
\subsection{Scope and caveats}\label{subsec: caveats}

\paragraph{Calibration identity versus existence of a calibrated model.} Propositions \ref{prop: calibration} and \ref{prop: main} are statements about a model that exists: \emph{if} (\ref{eq: threeFactorModel}) has a sufficiently regular solution, then it is calibrated if and only if its local volatility satisfies (\ref{eq: mainFormula}). Substituting (\ref{eq: mainFormula}) into (\ref{eq: threeFactorModel}) gives an SDE which is nonlinear in the sense of McKean, whose coefficient depends on the law of the solution through a conditional expectation given $S_t$. Neither existence nor uniqueness of a solution of this equation follows from the identity. The existence results above cover the stated Gaussian subclass and a regularized equation, not arbitrary unregularized smiles. Even for local-stochastic volatility models with deterministic rates, existence is known only in special cases: perturbative and local in time \cite{abergelTachet}, for regime-switching volatility under a condition on the range of the volatility factor \cite{jourdainZhou}, and for stationary solutions of a time-homogeneous version \cite{lackerShkolnikovZhang}. The more recent results of \cite{djete,mustapha} also concern different LSV structures and do not directly establish existence for the present general-smile stochastic-rates equation. The particle algorithm is therefore a numerical construction; its convergence as $N\to\infty$ is not proved in this paper.

\paragraph{Non-negativity of the local variance.}
The right hand side of (\ref{eq: mainFormula}) need not be positive.
The explicit Gaussian sufficient condition and the general necessary
log-contract bound in Section \ref{subsec: existence} have different scopes:
failure of a particular Gaussian construction is not a proof of nonexistence
among all spot-dependent models. In the numerical algorithm, report where
and by how much the variance floor binds. Such activation may reflect
estimation, interpolation or time-discretization error as well as incompatible
inputs; a negative trial iterate alone does not distinguish these cases.

\paragraph{Regularity.}
The hypotheses (H1)--(H3) distinguish local flow regularity and Malliavin
integrability from the density conditions needed for a classical strike-wise
formula. A global bound on $S\partial_S\sigma$ and a nonsingular correlation
matrix are not needed for smooth-payoff duality. In the algorithm, however,
noise in the estimated derivative $\partial_S\sigma_N$ enters the auxiliary
processes multiplicatively. Relaxing analytical assumptions does not remove
this numerical sensitivity.

\paragraph{Hull-White.} Proposition \ref{prop: main} uses two properties of the Hull-White model: the Clark-Ocone integrands (\ref{eq: kernels}) are deterministic and factorise, and the bond volatilities are deterministic, so that risk-neutral and forward measures differ by deterministic drifts.

\subsection{Outlook}
\paragraph{Existence.} In the system (\ref{eq: closureFive}), the equations of the two moments of $\ell^d$ and $\ell^f$ can be rewritten without second derivatives of the unknown variance $\sigma^2$ (a note in \texttt{paper/notes/} of the repository). Together with the known-denominator formulation of Theorem \ref{thm: closureNonlinear}, this is a possible starting point for a perturbative existence argument around the flat-smile solution of Proposition \ref{prop: gaussianExistence}. The estimates such an argument needs have not been obtained.

\paragraph{Gaussian multi-factor rates.} If a rate is a sum of Hull-White factors, the kernel (\ref{eq: kernels}) is a sum of exponentials, one per factor, and the matrix system (\ref{eq: matrixSystem}) extends with one Cameron-Martin direction per factor. The reduction to scalar processes has not been carried out for that case. Factors with different mean reversions cannot be aggregated into one process, and the cross sensitivities between each domestic and each foreign factor have to be tracked separately, so the number of auxiliary processes is not simply two per factor.

\paragraph{General one-factor short rates.} For rates $dr^i_t=\mu^i(t,r^i_t)dt+\sigma^i_r(t,r^i_t)dB^i_t$ the statement that replaces the factorisation $V_t/V_s$ is $\MD_sX_t=J_tJ_s^{-1}\mathcal{A}(s,X_s)\mathcal{C}^\top e$, with $J$ the $3\times3$ first variation matrix of $X=(S,r^d,r^f)$, which is not triangular because of the quanto drift. Two further difficulties appear: the Clark-Ocone integrand is stochastic and depends on $t$, and the forward measures differ from the risk-neutral ones by state-dependent drifts, which produce additional terms of the type of $H_t$.

\paragraph{Stochastic volatility.} With a stochastic volatility factor $a_t$ as in Section \ref{sec: oneRate}, the left hand side of (\ref{eq: mainFormula}) becomes $\sigma(t,K)^2\,\mathbb{E}^{d^t}[a_t^2|S_t=K]$ and the sensitivities acquire the tangent process of $a_t$ when $a_t$ is correlated with the rates. This is not worked out here.
